\section{Conclusion: the three answers}\label{sec:conclusion}
% claims used here: C6 C8 C15 C16 C24 C35 C37 m1-3 m2-8 m3-4

\paragraph{Question 1: what does $(A^k)_{ij}$ represent?}
For the adjacency matrix $A$ of a finite simple undirected graph and an integer
$k\ge1$, the entry $(A^k)_{ij}$ is the number of \emph{walks} of length $k$ from $n_i$ to
$n_j$ (\cref{thm:walkcount}). The proof is by induction on $k$, splitting the walks by their second-last
vertex. The number counts walks, not paths: a walk may repeat vertices and edges, so
$(A^2)_{ii}=\deg(n_i)>0$ although no path of length $2$ returns to $n_i$ (\cref{prop:warning}).
On the board graph, $(A^3)_{12}=2$ while the only path from $n_1$ to $n_2$ has length $1$.
The reading of the mark above $A$ as $k$ is an interpretation (\cref{sec:conventions}); the
theorem holds for any positive integer exponent.

\paragraph{Question 2: how can one tell from $A$ whether the graph is connected?}
For $N\ge2$, the graph is connected if and only if every entry off the diagonal of
$S=A+A^2+\dots+A^{N-1}$ is positive (\cref{thm:test}). Equivalently, $(I+A)^{N-1}$ has no zero
entry (\cref{prop:iplusa}). The proof uses walk counting, the lemma that every walk
contains a path, and the bound that a path has at most $N-1$ edges. The exponent $N-1$
cannot be lowered to $N-2$ (\cref{prop:sharp}). On the board graph the test says
connected, and on the graph with the single edge $12$ it says not connected, with the
components $\{1,2\}$ and $\{3\}$ (\cref{sec:example}). Breadth-first search gives the same
answers (\cref{sec:bfs}).

\paragraph{Question 3: what is a path between $n_i$ and $n_j$?}
In the course convention, a path between $n_i$ and $n_j$ is a walk $w_0,\dots,w_k$ with
$w_0=n_i$, $w_k=n_j$ and pairwise distinct vertices (\cref{def:pathbetween}). It is the
textbook's simple path, Definition~1.2.13, and the English path. The textbook's own word
``path'', Definition~1.2.12, means a walk with no repeated edge; this paper calls it a trail
(\cref{tab:conventions}). That the course uses this meaning comes from the assignment text and is
\kindtag{qualified}. Whether a path exists between two vertices does not depend on the choice
(\cref{prop:readings}).

\paragraph{What the paper did not settle.}
The board's mark above $A$, the extra edge and arrowhead in the board drawing, and the
lecturer's intended answer to Question~2 remain open (\cref{sec:open}). The statement that the textbook
contains no walk-counting theorem and no matrix test is \kindtag{qualified}: it rests on the pages checked and a full-text search (\cref{sec:provenance}).
